\section{The Angels in the Government of the World}\label{sec:government}

God governs all things, and he governs them in a manner worthy of his
goodness. Aquinas lays down the principle in the treatise on the divine
government. The plan of government (\latin{ratio gubernationis}), which is
providence itself, God holds immediately, down to the least particular; its
execution (\latin{executio}) he carries out in part through other things,
because it is a greater perfection for a thing to be good in itself and
also a cause of goodness in others. So God governs \latin{ut quasdam
aliarum in gubernando causas instituat, sicut si aliquis magister
discipulos suos non solum scientes faceret, sed etiam aliorum doctores}: as
a master who makes his pupils not only learned but teachers of others
(\ST{I q.~103 a.~6}). The angels stand first among these secondary
governors. Having shown how one angel acts on another by illumination,
speech, and hierarchy (qq.~106--109; \S\S\ref{sec:illumination},
\ref{sec:hierarchies}), Aquinas asks how the angels move bodies and how they
move men (\ST{I q.~106 pr.}). Question~110 treats \latin{de praesidentia
Angelorum super creaturam corporalem}, the presidency of the angels over the
corporeal creature; q.~111 their action on man by their natural power;
q.~112 their mission; q.~113 their guardianship (\ST{I q.~111 pr.};
\S\ref{sec:guardians}). The assault of the demons, the same natural powers
turned to harm by a perverse will, closes the treatise in q.~114
(\S\ref{sec:assaults}).

Scripture shows the angels at this work from end to end. They are
``mighty in strength, and execute his word, hearkening to the voice of his
orders,'' the ``ministers of his that do his will'' (Ps 102[103]:20--21).
The Apocalypse sees ``four angels standing on the four corners of the
earth, holding the four winds of the earth'' (Apoc 7:1), ``the angel of the
waters'' (Apoc 16:5), and an angel ``who had power over fire'' (Apoc
14:18). At the pond called Probatica ``an angel of the Lord descended at certain
times into the pond and the water was moved'' (John 5:2--4). The Law
was received ``by the disposition of angels'' (Acts 7:53), ``ordained by
angels in the hand of a mediator'' (Gal 3:19), and the Apostle gathers the
whole office into one question: ``Are they not all ministering spirits,
sent to minister for them who shall receive the inheritance of
salvation?'' (Heb 1:14; Douay throughout). The Fathers read these texts as
a doctrine of providence. The visible world is ordered through the
invisible; the angels serve that order as intelligent ministers; and
nothing in it escapes the will of God, who commands the good angels and
permits, restrains, and uses the wicked.

\sectionguard
\subsection{Augustine: The Corporeal Creation Ruled through the Spiritual}

The third book of Augustine's \emph{De Trinitate} is the patristic source
of the whole question, and Aquinas draws the \emph{sed contra} of three
articles of q.~110 from it. Augustine asks whether, in the appearances of
God to the fathers before the Incarnation, a creature was formed for the
purpose, or whether ``angels, who already existed, were so sent as to speak
in the person of God'' (\emph{De Trin.} III.1.4). To answer he sets out how
God governs bodies at all.

\paragraph{The will of God and the order of causes.} The principle stands
at the head: ``the power of the will of God reaches through the spiritual
creature even to visible and sensible effects of the corporeal creature''
(III.1.6). Augustine makes it plain with a wise man whose rational soul
consults the unchangeable truth and, at its bidding, wears out his body in
some work of mercy and falls ill. One physician blames dryness, another
moisture; both ``pronounce concerning proximate causes only, that is,
corporeal ones.'' Behind the illness lies the toil the soul undertook, and
behind the toil the command of wisdom, ``so nothing else but the will of
God would be found most truly to be the first cause of that illness''
(III.3.8). What holds of one wise man would hold of a city ruled by the
wise. Because it does not yet hold on earth, Augustine raises the mind to
the heavenly country. There the will of God, ``who maketh His angels
spirits, and His ministers a flaming fire,'' presides ``among spirits which
are joined in perfect peace and friendship, and combined in one will by a
kind of spiritual fire of charity,'' and ``thence diffuses itself through
all things by certain most perfectly ordered movements of the creature;
first spiritual, then corporeal'' (III.4.9). The sentence Aquinas
compresses into the \emph{sed contra} of \ST{I q.~110 a.~1} follows:

\begin{quote}
But as the more gross and inferior bodies are governed in due order by the
more subtle and powerful ones, so all bodies are governed by the living
spirit; and the living spirit devoid of reason, by the reasonable living
spirit; and the reasonable living spirit that makes default and sins, by
the living and reasonable spirit that is pious and just; and that by God
Himself \dots\ For nothing is done visibly or sensibly, unless either by
command or permission from the interior palace, invisible and
intelligible, of the supreme Governor, according to the unspeakable
justice of rewards and punishments, of favor and retribution, in that
far-reaching and boundless commonwealth of the whole creature.
(\emph{De Trin.} III.4.9)
\end{quote}

The Latin ladder runs \latin{ita omnia corpora per spiritum vitae, et
spiritus vitae irrationalis per spiritum vitae rationalem, et spiritus vitae
rationalis desertor atque peccator per spiritum vitae rationalem pium et
iustum}. The sinning spirit has its own rung. The demons stand inside the
order of government, ruled by the just spirits and through them by God,
and from the ``interior palace'' of the supreme Emperor, \latin{de
interiore invisibili atque intellegibili aula summi Imperatoris}, every
visible event is either commanded or permitted.

\paragraph{The usual course and the miracle.} The same divine power
``administers the whole spiritual and corporeal creature'' (III.5.11). God
draws rain from the sea every year; the rain that answered Elijah's prayer
after long drought was a miracle. God works ordinarily through thunder;
the thunder of Sinai, uttered so that ``certain intimations were given by
them,'' was a miracle. God draws the sap through the vine; the water made
wine at Cana was a miracle. When such things happen ``in a continuous kind
of river of ever-flowing succession \dots\ then they are called natural;
when, for the admonition of men, they are thrust in by an unusual
changeableness, then they are called miracles'' (III.6.11). Augustine names
the miracle by its rarity and its purpose of admonition. Aquinas will
define it more strictly by its cause (\ST{I q.~110 a.~4}).

\paragraph{The magicians and the finger of God.} Why then can magic arts do
the like, as when Pharaoh's wise men made serpents? Augustine turns the
question to the gnats, where the same magicians failed and confessed,
``This is the finger of God'' (Exod 8:19). ``And hence it is given us to
understand that not even those angels and powers of the air that
transgressed, who have been thrust down into that lowest darkness, as into
a peculiar prison, from their habitation in that lofty ethereal purity,
through whom magic arts have whatever power they have, can do anything
except by power given from above'' (III.7.12). That power is given for
three ends: to deceive the deceitful, as against the Egyptians; to warn the
faithful, lest they covet such things; and ``for the exercising, proving,
and manifesting of the patience of the righteous,'' as in Job.

\paragraph{The hidden seeds.} Here stands the \emph{sed contra} of
\ST{I q.~110 a.~2}: ``Yet it is not on this account to be thought that the
matter of visible things is subservient to the bidding of those wicked
angels; but rather to that of God'' (\latin{Nec ideo putandum est istis
transgressoribus angelis ad nutum servire hanc visibilium rerum materiam,
sed Deo potius}, III.8.13). The wicked angels through whom the magicians
made frogs and serpents did not create them. ``But, in truth, some hidden
seeds of all things that are born corporeally and visibly, are concealed in
the corporeal elements of this world'' (\latin{occulta quaedam semina}), and
the Creator of these invisible seeds is the Creator of all. As parents are
not the creators of men nor farmers of corn, ``so it is not right to think
not only the bad but even the good angels to be creators, if, through the
subtilty of their perception and body, they know the seeds of things which
to us are more hidden, and scatter them secretly through fit temperings of
the elements, and so furnish opportunities of producing things, and of
accelerating their increase.'' This is the \emph{sed contra} of \ST{I
q.~110 a.~3}. Augustine closes the paragraph with the law that governs all
angelic action on bodies:

\begin{quote}
But neither do the good angels do these things, except as far as God
commands, nor do the evil ones do them wrongfully, except as far as He
righteously permits. For the malignity of the wicked one makes his own will
wrongful; but the power to do so, he receives rightfully, whether for his
own punishment, or, in the case of others, for the punishment of the
wicked, or for the praise of the good. (\emph{De Trin.} III.8.13)
\end{quote}

Augustine's phrase ``perception and body'' belongs to his reserve about the
angels' bodies, on which he declines to decide whether they assume a
garment of the elements or change their own (III.1.5); Aquinas's account of
assumed bodies is in q.~51 (\S\ref{sec:q51}).

\paragraph{To create and to apply.} Augustine then draws the distinction on
which q.~110 turns. ``For it is one thing to make and administer the
creature from the innermost and highest turning-point of causation, which
He alone does who is God the Creator; but quite another thing to apply some
operation from without in proportion to the strength and faculties
assigned to each by Him'' (III.9.16). All things are already created ``in a
kind of texture of the elements,'' and ``as mothers are pregnant with
young, so the world itself is pregnant with the causes of things that are
born.'' The angels, by the subtlety of their perception and the speed of
their motion, apply these causes more quickly than men can, and the
quickness is itself a wonder to us (III.9.17). Their power has a double
limit, and Augustine states it in an image:

\begin{quote}
For we know that a man can walk, yet that he cannot do so if he is not
permitted; but that he cannot fly, even if he be permitted. So those
angels, also, are able to do certain things if they are permitted by more
powerful angels, according to the supreme commandment of God; but cannot
do certain other things, not even if they are permitted by them; because
He does not permit from whom they have received such and such a measure of
natural powers. (\emph{De Trin.} III.9.18)
\end{quote}

Some things the wicked angels could do by nature and cannot because they
are not permitted; other things lie beyond their nature altogether. Aquinas
takes the double limit into \emph{De potentia}, where he explains that the
demons are said to be bound (\latin{ligari}) when they are hindered from
what their nature could reach, and loosed when God's judgment permits it
(\emph{De pot.} q.~6 a.~5).

\paragraph{The theophanies wrought by angels.} Beyond the usual course of
nature there are signs ``brought within reach of our senses in order to
announce something from God'' (III.10.19). The deed of an angel and the
deed of a man may signify the same thing, and yet ``there is a wide
difference between the deed of the angel and the deed of the man. The
former is both to be wondered at and to be understood, the latter only to
be understood.''
Augustine likens it to ``the name of God \dots\ written both in gold and in
ink'': the gold is more precious, the thing signified one (III.10.20). He
confesses the limit of his sight: ``how the angels do these things, or
rather, how God does these things by His angels, and how far He wills them
to be done even by the bad angels, whether by permitting, or commanding, or
compelling, from the hidden seat of His own supreme power; this I can
neither penetrate by the sight of the eyes, nor make clear by assurance of
reason, nor be carried on to comprehend it by reach of intellect''
(III.10.21). Aquinas sets these words at the head of his disputed question
on the powers of spiritual creatures and proceeds, as Augustine taught,
\latin{absque assertione et sententiae melioris praeiudicio} (\emph{De
pot.} q.~6 a.~3). What Augustine does affirm he affirms from Scripture:
``all those appearances to the fathers \dots\ were wrought through the
creature. And if we cannot discern in what manner He wrought them by
ministry of angels, yet we say that they were wrought by angels''
(III.11.22). The proof is Hebrews, where ``the word spoken by angels''
(Heb 2:2) is set against the word spoken by the Son; Stephen, who says that
there appeared to Moses ``an angel in a flame of fire in a bush'' and that
the voice said, ``I am the God of thy fathers'' (Acts 7:30--32); the angel who called to
Abraham on the mountain and said ``for my sake'' in the person of God (Gen
22:11--12); and the Law ``received \dots\ by the disposition of angels''
(Acts 7:53; Gal 3:19). Why then is it written ``The Lord said to Moses''?
``Because, when the crier proclaims the words of the judge, it is not
usually written in the record, so and so the crier said, but so and so the
judge'' (III.11.23). The book ends on its conclusion: the words and bodily
appearances given to the fathers before the Incarnation, ``when God was
said to appear, were wrought by angels'' (III.11.27).

\paragraph{Contemplating and commanding at once.} In the \emph{De Genesi ad
litteram} Augustine states the same order as the angels' own share in
providence. To the exalted angels, who enjoy God and serve him in
blessedness, is subject \latin{omnis natura corporea, omnis irrationalis
vita, omnis voluntas vel infirma vel prava}, so that they do with their
subjects what the order of nature requires, at the command of him to whom
all things are subject. They see the unchangeable truth in God and direct
their wills by it; they are moved in time by his command, \latin{nec ita
ut ab eius contemplatione resiliant aut defluant; sed simul et illum sine
loco ac tempore contemplantur, et eius in inferioribus iussa perficiunt}:
they contemplate him and carry out his orders in lower things at once, and
do not fall back from contemplation in doing so (\emph{De Gen.\ ad litt.}
VIII.24.45). In the \emph{Eighty-Three Questions} he gives the rule Aquinas
quotes for the holy doctors: \latin{unaquaeque res visibilis in hoc mundo
habet potestatem angelicam sibi praepositam}, every visible thing in this
world has an angelic power set over it (\emph{De div.\ qq.\ 83} q.~79.1).

\sectionguard
\subsection{How Angels Act on Bodies (\ST{I q.~110})}\label{sec:q110}

Four articles: whether the corporeal creature is governed by the angels;
whether corporeal matter obeys the mere will of the angels; whether the
angels by their own power can immediately move bodies locally; and whether
the good or bad angels can work miracles (\ST{I q.~110 pr.}).

\paragraph{The corporeal creature is governed by the angels (a.~1).} The
\emph{sed contra} joins two Latin Doctors. Augustine says ``that `all
bodies are ruled by the rational spirit of life''\,'' (\emph{De Trin.}
III.4, as Aquinas compresses Augustine's ladder of causes), and Gregory ``that `in this
visible world nothing takes place without the agency of the invisible
creature'\,'' (\emph{Dial.} IV). Aquinas argues from a law common to human
and natural affairs: ``every particular power is governed and ruled by the
universal power; as, for example, the bailiff's power is governed by the
power of the king.'' The superior angels, who have more universal
knowledge, rule the inferior. The power of any body is more particular than
that of any spirit, ``for every corporeal form is a form individualized by
matter, and determined to the `here and now'; whereas immaterial forms are
absolute and intelligible.'' Therefore ``as the inferior angels who have the
less universal forms, are ruled by the superior; so are all corporeal
things ruled by the angels. This is not only laid down by the holy doctors,
but also by all philosophers who admit the existence of incorporeal
substances.'' A body has determinate actions but exercises them only when
moved, ``because it belongs to a body not to act unless moved'' (ad~1).
Aristotle gave the separate substances the heavenly motions alone, since
he saw nothing in the lower world beyond what the heavens suffice to
cause; but ``because we assert that many things are done in the inferior
bodies besides the natural corporeal actions,'' Aquinas concludes ``that
the angels possess an immediate presidency not only over the heavenly
bodies, but also over the inferior bodies'' (ad~2). The third reply sets
the philosophers side by side. Plato set immaterial substances as the
types of sensible things, each immediately over its own; Aristotle gave
them only the universal agents, the heavens; Avicenna took a middle way,
with one ``active intelligence'' over all lower bodies. ``The holy doctors
held with the Platonists that different spiritual substances were placed
over corporeal things'': Augustine in the \emph{Eighty-Three Questions},
Damascene on the devil as one of the powers ``who presided over the
terrestrial order,'' and Origen on Balaam's ass (Num 22:23), that the world
needs angels ``who preside over beasts, and over the birth of animals, and
trees, and plants'' (\emph{Hom.\ in Num.} 14, as Aquinas cites it). The
reason is not that one angel is fitted by nature to animals and another to
plants, ``because each angel, even the least, has a higher and more
universal power than any kind of corporeal things: the reason is to be
sought in the order of Divine wisdom, Who places different rulers over
different things.'' These charges do not multiply the orders beyond nine,
for orders are distinguished by general offices; as the angels set over
the demons belong to the Powers, ``so to the order of the `virtues' do those
angels seem to belong who preside over purely corporeal creatures; for by
their ministration miracles are sometimes performed'' (ad~3).

Gregory's sentence belongs to his proof in the \emph{Dialogues} that the
soul lives on after death. The body moves only while the soul is in it;
the eye sees only by the invisible soul; the builder's machines lift
timber and stone only because an invisible life moves the visible body.
``By which we may easily gather, that nothing can be disposed of in this
visible world, but by another creature which is invisible: for as almighty
God either by inspiration, or by replenishing those creatures which have
reason, doth both quicken and move those things which be invisible, so, in
like manner, those things which be invisible do give motion and sense to
carnal bodies which are visible'' (\emph{Dial.} IV.5).\footnote{Gregory,
\emph{Dialogues} IV, English of 1608 as edited by E.~G. Gardner (1911),
where the passage stands in chapter~5; Aquinas cites it as \emph{Dial.}
IV.6.} Damascene teaches the same presidency from the East. The angels ``are
the guardians of the divisions of the earth: they are set over nations and
regions, allotted to them by their Creator: they govern all our affairs and
bring us succour,'' and the reason is ``because they are set over us by the
divine will and command and are ever in the vicinity of God'' (\emph{De
fide orth.} II.3). The chief of the fallen, before his fall, was the one
``set over the earthly realm, and into whose hands God committed the
guardianship of the earth'' (II.4).

The \emph{Summa contra gentiles} gives the article its fullest reasons.
Because providence descends from the highest to the lowest in proportion,
and the intellectual creatures are highest, \latin{exigit igitur divinae
providentiae ratio ut ceterae creaturae per creaturas rationales regantur}
(\emph{SCG} III.78 n.~1). The intellectual creature shares both parts of
providence, the ordering that belongs to knowledge and the execution that
belongs to operative power, where other creatures share only execution
(n.~2); it alone knows the reasons of the order and acts of itself by free
will (nn.~5--6). The same chapter series assigns the offices of the orders
in the government of the world, and there the Virtues move the heavens and
carry out \latin{divinorum operum quae praeter naturae ordinem fiunt},
since Gregory says \latin{virtutes dicuntur illi spiritus per quos signa
frequentius fiunt} (\emph{SCG} III.80 n.~11; Gregory, \emph{Hom.\ in Ev.}
34.10).

\paragraph{Matter does not obey the mere will of an angel (a.~2).} The
\emph{sed contra} is Augustine's sentence from \emph{De Trinitate} III.8.13:
``It is not to be thought, that this visible matter obeys these rebel
angels; for it obeys God alone.'' The adversaries are the Platonists, who
made material forms participations of immaterial ones, and Avicenna, who
held that all forms in matter proceed from the concept of an intelligence
while bodies only dispose matter. Aquinas answers from Aristotle: what is
made, properly speaking, is the composite, for the form is ``not as that
which is, but as that by which something is.'' Every agent makes its like,
so what makes natural things either is itself composite, ``as when fire
begets fire,'' or has the whole composite within its power, which belongs
to God alone. Hence \latin{omnis informatio materiae vel est a Deo
immediate, vel ab aliquo agente corporali; non autem immediate ab
Angelo}: every informing of matter is immediately from God or from some
corporeal agent, and never immediately from an angel. Our soul changes its
own body by a concept because it is that body's form, and the angel has no
such union with natural bodies (ad~1). The superior power does what the
inferior does ``in a more excellent way,'' by moving the corporeal agents
themselves (ad~2). A natural effect may follow from angelic power that no
bodily agent could produce alone, yet matter does not thereby obey an
angel's will, ``as neither does matter obey the mere will of a cook, when
by regulating the fire according to the prescription of his art he
produces a dish that the fire could not have produced by itself'' (ad~3).

The article is Augustine's doctrine of the seeds and Damascene's rule
together. ``Those who say that the angels are creators of any kind of
essence whatever are the mouth of their father, the devil. For since they
are created things they are not creators'' (\emph{De fide orth.} II.3).
\emph{De potentia} names the same adversary and adds that the faith agrees
with Aristotle here, \latin{unde Augustinus dicit, quod Angelis non servit
ad nutum materia corporalis} (\emph{De pot.} q.~6 a.~3); \emph{De malo}
concludes that spiritual substances \latin{propria virtute non possunt
formaliter transmutare inferiora corpora, nisi adhibendo aliqua corporalia
activa proportionata effectibus quos intendunt; sicut homo potest calefacere
per ignem} (\emph{De malo} q.~16 a.~9).

\paragraph{Bodies obey the angels in local motion (a.~3).} The \emph{sed
contra} is again Augustine: the angels ``use corporeal seed to produce
certain effects'' (\emph{De Trin.} III.8--9), which they cannot do without
local motion. Aquinas argues from Dionysius: ``Divine wisdom has joined the
ends of the first to the principles of the second'' (\emph{DN} 7). In
Parker's version the divine wisdom is ``ever uniting the ends of the former
to the beginnings of those that follow, and beautifying the one symphony
and harmony of the whole'' (\emph{DN} 7.3). The lower nature touches the
higher at its summit. Corporeal nature lies below spiritual, and the most
perfect of bodily movements is local motion, because ``what is moved
locally is not as such in potentiality to anything intrinsic, but only to
something extrinsic---that is, to place.'' Therefore bodies have a natural
aptitude to be moved immediately by spirits as to place, as the soul moves
its body ``first and chiefly by a local motion.'' Local motions need not
follow from a body's form, as the ebb and flow of the sea follows the moon
(ad~1). By local motion, the first of motions, the angels cause the
others, ``employing corporeal agents to produce these effects, as a workman
employs fire to soften iron'' (ad~2). The soul's motive power is confined to
its own body; ``an angel's power is not limited to any body; hence it can
move locally bodies not joined to it'' (ad~3). \emph{De malo} adds a second
reason: local motion is the least change a body can undergo, touching only
its place, and the weakest effects can come immediately from the most
remote agent (\emph{De malo} q.~16 a.~10). \emph{De potentia} calls the
resulting works a kind of art. The angels move by their command bodies in
which there is a natural active power, the \latin{naturae semina} of
Augustine, \latin{et sic operatio eorum non erit per modum miraculi, sed
per modum artis}; and their art surpasses man's because they know the
powers of bodies better and apply them faster (\emph{De pot.} q.~6 a.~3).

\paragraph{God alone works miracles (a.~4).} The \emph{sed contra} is Ps
135[136]:4, of God: ``Who alone doth great wonders.'' The first objection is
Gregory's sentence on the Virtues, ``by whom signs and miracles are usually
done'' (\emph{Hom.\ in Ev.} 34.10). Aquinas defines: ``A miracle properly so
called is when something is done outside the order of nature,'' but not
outside the order of some particular nature, ``for otherwise anyone would
perform a miracle by throwing a stone upwards.'' A miracle is what is done
\latin{praeter ordinem totius naturae creatae}, outside the order of the
whole created nature. ``But God alone can do this, because, whatever an
angel or any other creature does by its own power, is according to the
order of created nature; and thus it is not a miracle. Hence God alone can
work miracles.'' The angels are said to work miracles in two ways: ``either
because God works miracles at their request, in the same way as holy men
are said to work miracles; or because they exercise a kind of ministry in
the miracles which take place; as in collecting the dust in the general
resurrection'' (ad~1). What a creature does by a power unknown to us is a
miracle ``as regards ourselves,'' and in this sense magicians work wonders
through the demons (ad~2). Aquinas uses Augustine's three-fold distinction
from the \emph{Eighty-Three Questions}: ``magicians work miracles by
private contracts; good Christians by public justice, bad Christians by the
signs of public justice.'' Every power of a creature is like a private
person in a city, while ``the Divine justice is in the whole universe as
the public law is in the city'' (ad~2). At its own locus Augustine's
sentence reads \latin{magi per privatos contractus, boni christiani per
publicam iustitiam, mali christiani per signa publicae iustitiae}, and he
adds that the powers dare not despise the signs of the Emperor:
\latin{contremiscunt enim haec, ubicumque illa conspexerint} (\emph{De div.\
qq.\ 83} q.~79.4). The angels can act outside some particular order of
bodily nature, but never outside the whole created order (ad~4). By this
definition Augustine's broader name, the miracle as a work thrust in by an
unusual change ``for the admonition of men'' (\emph{De Trin.} III.6.11), is
kept for what is wonderful to us, and the strict name is reserved to God.

The parallels press the article further. \emph{Contra gentiles} answers
Avicenna, who held that matter obeys the mere apprehension of separate
substances, and concludes \latin{relinquitur quod substantiae spirituales
creatae non faciant miracula propria virtute}; yet \latin{nihil prohibet
huiusmodi substantias, inquantum agunt in virtute divina, miracula
facere}, which is why Gregory assigns one order specially to miracles
(\emph{SCG} III.103 nn.~6--7). \emph{De potentia} names three ways in which
the angels share in miracles: by prayer, as men do; by disposing matter
with their natural power, as in gathering the dust of the dead; and by
cooperating as instruments of the divine command, as a grace given for the
moment (\emph{De pot.} q.~6 a.~4).

\angeldossier{\ST{I q.~110 aa.~1--4}; parallels \emph{Summa contra
gentiles} III.78, 80, 103; \emph{De potentia} q.~6 aa.~3--5; \emph{De malo}
q.~16 aa.~9--10.}
 {Church teaching that the devil creates nothing and does not produce
 thunder, lightning, storms, or droughts by his own authority (Braga~I,
 can.~8, DS~458), and that Satan's power is a creature's (CCC~395). Common
 teaching that God governs the corporeal creation through the angels, the
 wicked included under his command or permission (a.~1; Augustine,
 \emph{De Trin.} III.4.9; Gregory, \emph{Dial.} IV; Damascene, \emph{De
 fide orth.} II.3; Ps 102[103]:20--21), and that angels are not creators
 and matter does not serve their bare will (a.~2; Augustine, \emph{De
 Trin.} III.8.13; Damascene II.3). Thomist position that every informing
 of matter is immediately from God or from a corporeal agent (a.~2), that
 bodies obey the angels immediately in local motion and in other changes
 through it (a.~3), and that the Virtues preside over purely corporeal
 things (a.~1 ad~3). Common teaching that miracles in the proper sense are
 God's alone and that angels share in them by prayer and ministry (a.~4;
 Ps 135[136]:4).}
 {Augustine, \emph{De Trin.} III.4.9, III.8.13, III.9.16--18 (\emph{sed
 contra} of aa.~1--3), \emph{De div.\ qq.\ 83} q.~79 (a.~1 ad~3; a.~4
 obj.~2, ad~2), \emph{De Gen.\ ad litt.} VIII.24.45; Gregory, \emph{Dial.}
 IV (\emph{sed contra} of a.~1) and \emph{Hom.\ in Ev.} 34.10 (a.~4
 obj.~1), read at the locus; Damascene, \emph{De fide orth.} II.3--4 (a.~1
 ad~3), read at the locus; Origen, \emph{Hom.\ in Num.} 14 (a.~1 ad~3), cited by
 Aquinas; sentences cited under Augustine's name in a.~4 obj.~2--3, which the 1920
 English edition assigns to the pseudo-Augustinian \emph{Liber XXI
 sententiarum}; Dionysius, \emph{DN} 7.3 (a.~3); Aristotle,
 \emph{Metaph.} VII, \emph{Phys.} VIII; Plato and Avicenna as reported;
 Ps 135[136]:4, Num 22:23, Exod 7--8; Braga~I (DS~458); CCC~395.}
 {Plato, who set immaterial forms immediately over each kind of body, and
 Avicenna, for whom the forms of matter flow from the active intelligence
 and matter obeys the apprehension of separate substances (aa.~1--2;
 \emph{SCG} III.103). Aristotle, who gave the separate substances no
 immediate rule over lower bodies (a.~1 ad~2), and Alexander, who referred
 to the heavenly bodies every effect ascribed to angels and demons
 (\emph{De pot.} q.~6 a.~3). Augustine names miracles by their rarity and
 purpose of admonition (\emph{De Trin.} III.6.11), where Aquinas defines
 them by their exceeding the whole created order (a.~4).}

\sectionguard
\subsection{The Action of the Angels on Man (\ST{I q.~111})}\label{sec:q111}

Four articles, on what an angel can do to a man by his natural power:
whether he can enlighten the human intellect; whether he can change the
will; whether he can change the imagination; and whether he can change the
senses (\ST{I q.~111 pr.}). The answers descend through the powers of the
soul, and the good and the wicked angel are treated together as to natural
power. The difference lies in the use: the good angel serves God's
enlightening of man, the wicked angel turns the same powers against it
(\S\ref{sec:q114}).

\paragraph{The angels enlighten man (a.~1).} The \emph{sed contra} is
Dionysius, who proves ``that the revelation of Divine things reaches men
through the ministry of the angels'' (\emph{CH} 4). At its locus the fourth
chapter explains why the name angel belongs pre-eminently to the heavenly
beings: ``the supremely Divine illumination comes to them at first hand,
and, through them, there pass to us manifestations above us. Thus, then,
the Law, as the Word of God affirms, was given to us through the
ministration of Angels; and Angels led our illustrious fathers before the
Law, and after the Law, to the Divine Being'' (\emph{CH} 4.2). The fathers
``were initiated into these Divine visions, through the mediation of the
Heavenly Powers'' (4.3), and even the mystery of Christ's love for man
passed first to the angels and through them to us: Gabriel to Zachary and
to Mary, another angel to Joseph, another to the shepherds (4.4). In the
\emph{Divine Names} the angels are ``as it were, heralds of the Divine
silence,'' who ``project, as it were, luminous lights revealing Him Who is
in secret,'' and the souls of men are able ``to be conducted by them as
good guides to the good Origin of all good things'' (\emph{DN} 4.2).

Aquinas argues from the order of providence: ``as the inferior angels are
enlightened by the superior, so men, who are inferior to the angels, are
enlightened by them.'' Illumination has two parts, as Aquinas showed of
the angels themselves (\ST{I q.~106 a.~1}; \S\ref{sec:q106}): the inferior intellect is
strengthened by the superior, and the intelligible species in the superior
are proposed to the inferior so that it can grasp them. A superior angel
divides his universal concept for the lower angel. Man cannot grasp the
intelligible truth unveiled, since he understands by turning to phantasms,
``so the angels propose the intelligible truth to men under the
similitudes of sensible things,'' according to Dionysius: ``It is
impossible for the divine ray to shine on us, otherwise than shrouded by
the variety of the sacred veils'' (\emph{CH} 1). Parker's version reads:
``it is not possible that the supremely Divine Ray should otherwise
illuminate us, except so far as it is enveloped, for the purpose of
instruction, in variegated sacred veils'' (\emph{CH} 1.2). ``On the other
hand, the human intellect as the inferior, is strengthened by the action
of the angelic intellect'' (\latin{fortificatur per actionem intellectus
angelici}).

The replies place the angels inside the act of faith without making them
its author. Faith has two parts. The inclination of the intellect to
assent at the will's command is from God alone, since ``the intellect
assents to the truth of faith, not as convinced by the reason, but as
commanded by the will; hence Augustine says, `No one believes except
willingly.'\,'' But what is to be believed must be proposed, ``which is
accomplished by man, according to Romans 10:17, `Faith cometh by hearing';
principally, however, by the angels, by whom Divine things are revealed to
men. Hence the angels have some part in the enlightenment of faith,'' and
not only in what is to be believed ``but also as regards what is to be
done'' (ad~1). Natural reason, which is immediately from God, can be
strengthened by an angel, ``so that he may obtain from creatures a more
perfect knowledge of God'' (ad~2). A man knows that he understands, but
``not everyone who is enlightened by an angel, knows that he is enlightened
by him'' (ad~3). \emph{De malo} states the difference between the good angel
and the wicked here. As to the light, the angelic light can strengthen the
human, \latin{quod Angelus bonus intendit, non autem Angelus malus. Unde
hoc modo Angeli boni movent animam ad intelligendum, non autem Daemones}
(\emph{De malo} q.~16 a.~12). As to the species, both can arrange the
images of the imagination, the good for man's good, the demons for his
harm, either to lead him into pride or to entangle him in doubts he cannot
resolve (ibid.).

\paragraph{The will is changed by God alone (a.~2).} The \emph{sed contra}
is Prov 21:1: ``the heart of the king is in the hand of the Lord:
whithersoever he will, he shall turn it.'' The will can be changed from
within, and so only by God, ``because He gives the power of such an
inclination to the intellectual nature.'' It can be moved from without by
the good apprehended, and so ``God alone can move the will efficaciously;
but an angel and man move the will by way of persuasion'' (\latin{per
modum suadentis}). In man there is a third way, through the passions of
the sensitive appetite, and here ``the angels, as being able to rouse these
passions, can move the will, not however by necessity, for the will ever
remains free to consent to, or to resist, the passion.'' The angels ``burn
away vices'' by persuasion, as God's ministers (ad~1). The second objection
set Bede, who says the devil ``does not send wicked thoughts, but kindles
them,'' against Damascene, who says the evils are ``contrived by the demons
and put into men'' (\emph{De fide orth.} II.4, as Aquinas renders it).
Aquinas reconciles them: ``The demon cannot put thoughts in our minds by
causing them from within, since the act of the cogitative faculty is
subject to the will; nevertheless the devil is called the kindler of
thoughts, inasmuch as he incites to thought, by the desire of the things
thought of, by way of persuasion, or by rousing the passions. Damascene
calls this kindling `a putting in' because such a work is accomplished
within. But good thoughts are attributed to a higher principle, namely,
God, though they may be procured by the ministry of the angels'' (ad~2).
The imagination can be changed by an angel, but the will can follow reason
against the passions (ad~3).

The Fathers had taught the article before it was posed. Origen holds that
thoughts ``sometimes proceed from ourselves, and sometimes are originated
by the opposing powers; not seldom also are they suggested by God, or by
the holy angels,'' and he adds the limit: ``We are not, however, to imagine
that any other result follows from what is suggested to our heart, whether
good or bad, save a (mental) commotion only, and an incitement instigating
us either to good or evil'' (\emph{De princ.} III.2.4). Damascene says of
the demons: ``while the liberty to attack man has been granted to them,
they have not the strength to over-master any one: for we have it in our
power to receive or not to receive the attack'' (\emph{De fide orth.}
II.4). Abbot Serenus in Cassian teaches that the enemy opposes our progress
``in such a way that we can think of them as only inciting to evil things
and not forcing'' (\emph{Conl.} VII.8), and that spirit cannot be implanted
in spirit so as to hold it, ``for this is the true prerogative of Deity
alone'' (VII.10). \emph{De malo} draws the consequence for knowledge: the
act of the will is known to God and to the one who wills, and to no other
(\emph{De malo} q.~16 a.~8; \S\ref{sec:q57}).

\paragraph{The imagination (a.~3).} The \emph{sed contra} is the angel who
appeared to Joseph ``in his sleep'' (Matt 1:20; 2:13). ``Both a good and a
bad angel by their own natural power can move the human imagination.''
Bodies obey the angels in local motion (\ST{I q.~110 a.~3}), and
imaginative apparitions are sometimes caused ``by the local movement of
animal spirits and humors,'' as Aristotle explains dreams, the impressions
of past sensation moving the sensitive principle in sleep. What happens by
a natural disturbance of the humors, and by a man's voluntary imagining,
can be done by the power of an angel, ``sometimes with alienation from the
bodily senses, sometimes without such alienation.'' The angel does not
impress a form never received from sense, ``for he cannot make a man born
blind imagine color'' (ad~2). Augustine speaks of spirit mingling with
spirit (\emph{De Gen.\ ad litt.} XII, as Aquinas cites it); the
commingling ``is not a mingling of essences, but by reason of an effect
which he produces in the imagination \dots\ so that he shows man what he
[the angel] knows, but not in the way he knows'' (ad~3). When a good angel
causes an imaginative vision he sometimes enlightens the intellect as
well, and then there is no deception; when images only appear, any
deception comes from the defect of the one who sees, ``thus neither was
Christ a cause of deception when He spoke many things to the people in
parables, which He did not explain to them'' (ad~4). \emph{De malo} gives
the same account at length and rejects the alternatives Augustine had
weighed: the soul does not see the essence of a spirit or the species in
it, and an angel cannot infuse new forms into sense or imagination, so the
change must come from the local movement of what already exists in the
body (\emph{De malo} q.~16 a.~11).

\paragraph{The senses (a.~4).} The \emph{sed contra} is the blindness with
which the angels struck the men of Sodom, ``so that they could not find the
door'' (Gen 19:11), and the Syrians whom Eliseus led blinded into Samaria
(4 Kgs 6:18). The senses are changed from without by a sensible object and
from within by the disturbance of spirits and humors, ``as for example, a
sick man's tongue, charged with choleric humor, tastes everything as
bitter.'' An angel can work both: he ``can offer the senses a sensible
object from without, formed by nature or by the angel himself, as when he
assumes a body'' (\S\ref{sec:q51}), and he can move spirits and humors
within. The interior principle of sensation remains the soul's own, but it
can be moved from without (ad~1). By the same movement of spirits an angel
can affect the nutritive and appetitive powers and any power using a bodily
organ (ad~2). ``An angel can do nothing outside the entire order of
creatures; but he can outside some particular order of nature,'' and so
``in some special way an angel can work a change in the senses outside the
common mode of nature'' (ad~3).

\angeldossier{\ST{I q.~111 aa.~1--4}; parallels \emph{De malo} q.~16
aa.~8, 11--12; \emph{SCG} III.81.}
 {Common teaching that the angels enlighten men by proposing divine truth
 and strengthening the mind, and that revelation, the Law included, came
 through their ministry (a.~1; Acts 7:53; Gal 3:19; Dionysius, \emph{CH}
 4.2--4). Common teaching that no angel or demon can move the will from
 within or compel it, God alone moving it efficaciously; angels persuade
 and can rouse passions, and the will remains free to consent or resist
 (a.~2; Prov 21:1; Origen, \emph{De princ.} III.2.4; Damascene, \emph{De
 fide orth.} II.4; Cassian, \emph{Conl.} VII.8, 10). Common teaching that
 good and wicked angels can act on imagination and senses (aa.~3--4; Matt
 1:20; Gen 19:11; 4 Kgs 6:18). Thomist position on the manner: species
 proposed under sensible likenesses, and imagination and sense moved
 through the local motion of spirits and humors, never by an infused form
 (aa.~1, 3--4; \emph{De malo} q.~16 a.~11).}
 {Dionysius, \emph{CH} 4 (\emph{sed contra} of a.~1) and \emph{CH} 1.2
 (a.~1), read at the locus in Parker; \emph{EH} on baptism as
 illumination (a.~1 obj.~1) and \emph{DN} 4.1--2, the latter read at the
 locus; Augustine on belief and the will (a.~1 ad~1), \emph{De Gen.\ ad
 litt.} XII.12 (a.~3 obj.~3), cited by Aquinas; Bede on Matt 15 (a.~2
 obj.~2), cited by Aquinas; Damascene, \emph{De fide orth.} II.4, read at
 the locus; Origen, \emph{De princ.} III.2.4, and Cassian, \emph{Conl.}
 VII.8--10, read at the locus; Aristotle, \emph{De somno}; the gloss on
 Heb 1:7 and Rom 1:19; Prov 21:1, Rom 10:17, Eph 2:8, Matt 1:20, 2:13,
 Gen 19:11, 4 Kgs 6:18.}
 {Damascene's sentence that the demons put evils into men, read by Aquinas
 as a kindling from within the soul's powers and not a causing of thought
 from within the mind (a.~2 ad~2). Augustine's three possibilities for
 imaginative visions (\emph{De Gen.\ ad litt.} XII), of which \emph{De
 malo} q.~16 a.~11 admits only the movement of what pre-exists in the
 body.}

\sectionguard
\subsection{The Mission of the Angels (\ST{I q.~112})}\label{sec:q112}

Four articles: whether any angels are sent on works of ministry; whether
all are sent; whether those who are sent also assist; and from which
orders they are sent (\ST{I q.~112 pr.}). The question gathers the
Dionysian law that the highest are not sent, Gregory's distinction of the
assisting and the ministering angels from Dan 7:10, and the reconciliation
of both with the Gospel's word that the angels of the little ones ``always
see the face of my Father'' (Matt 18:10).

\paragraph{The angels are sent in ministry (a.~1).} The \emph{sed contra}
is Exod 23:20: ``Behold I will send my angel, who shall go before thee.''
Aquinas takes the notion of mission from the treatise on the Trinity: ``he
is said to be sent who in any way proceeds from another so as to begin to
be where he was not, or to be in another way, where he already was.'' The
Son and the Spirit are sent as proceeding by origin and beginning to be in
a new way, by grace or by the nature assumed, where they were already by
the presence of Godhead. God is everywhere as universal agent; an angel's
power, as a particular agent, ``reaches to one thing in such a way as not to
reach another; and so he is `here' in such a manner as not to be `there.'\,''
When an angel is to work on some body, he applies his power to it anew and
so begins to be there, and all this at God's command. The action ``proceeds
from God as from its first principle, at Whose nod and by Whose authority
the angels work; and is reduced to God as to its last end. Now this is what
is meant by a minister: for a minister is an intelligent instrument''
(\latin{minister est sicut instrumentum intelligens}). ``Hence angels'
actions are called `ministries'; and for this reason they are said to be
sent in ministry.'' Damascene describes the same circumscription: ``They are
circumscribed: for when they are in the Heaven they are not on the earth:
and when they are sent by God down to the earth they do not remain in the
Heaven,'' and ``their nature is endowed with such celerity that wherever
the Divine glance bids them there they are straightway found'' (\emph{De
fide orth.} II.3).

The replies meet four difficulties. An intellectual act that dwells in the
intellect, as contemplation, occupies no place, but an act commanded by
intellect may have a place (ad~1). The empyrean heaven befits the angels'
dignity by congruity; an angel absent from it loses nothing, ``as neither
does a king lessen his dignity when not actually sitting on his regal
throne'' (ad~2). Outward occupation dims our contemplation because we act
through the senses; ``an angel, on the contrary, regulates his exterior
actions by intellectual operation alone. Hence it follows that his external
occupations in no respect impede his contemplation; because given two
actions, one of which is the rule and the reason of the other, one does not
hinder but helps the other,'' as Gregory says, ``the angels do not go abroad
in such a manner as to lose the delights of inward contemplation'' (ad~3).
The angels minister chiefly to God and only secondarily to us, ``not
because we are superior to them, absolutely speaking, but because, since
every man or angel by cleaving to God is made one spirit with God, he is
thereby superior to every creature'' (ad~4; Phil 2:3).

Gregory's sentence stands in the \emph{Moralia} on ``the sons of God'' who
``came to stand before the Lord'' (Job 1:6). He asks from where the angels
can come, when they ``always see the face of my Father'' (Matt 18:10) and
``stood before him'' (Dan 7:10), and yet Paul calls them ``ministering
spirits, sent to minister'' (Heb 1:14). His answer lies in the subtlety of the angelic nature: ``For
they never so go forth apart from the vision of God, as to be deprived of
the joys of interior contemplation; for if when they went forth they lost
the vision of the Creator, they could neither have raised up the fallen,
nor announced the truth to those in ignorance.'' The angels are bounded by
place, yet their knowledge reaches beyond ours ``since they contemplate the
very source of knowledge itself.'' ``Therefore they are both sent from Him,
and stand by Him too, since both in that they are circumscribed, they go
forth, and in this, that they are also entirely present, they never go
away. Thus they at the same time always behold the Father's face, and yet
come to us'' (\emph{Moralia} II.3.3). The homily on the angels repeats it
in the form Aquinas quotes: \latin{Et mittuntur igitur, et assistunt, quia
etsi circumscriptus est angelicus spiritus, summus tamen spiritus ipse,
qui Deus est, circumscriptus non est. Angeli itaque et missi, et ante ipsum
sunt, quia quolibet missi veniant, intra ipsum currunt}: wherever they are
sent, they run within God (\emph{Hom.\ in Ev.} 34.13). Augustine had said
the same of the angels who contemplate God without place and time and at
once carry out his commands below (\emph{De Gen.\ ad litt.} VIII.24.45).

\paragraph{The higher angels are not sent to outward ministry (a.~2).} The
\emph{sed contra} is Gregory ``quoting the statement of Dionysius'' that
``the higher ranks fulfil no exterior service'' (\latin{superiora agmina
usum exterioris ministerii nequaquam habent}). At its locus Gregory calls
Dionysius \latin{antiquus videlicet et venerabilis Pater} and reports his
teaching \latin{quod ex minoribus angelorum agminibus foras ad explendum
ministerium vel visibiliter vel invisibiliter mittuntur, scilicet quia ad
humana solatia ut angeli aut archangeli veniunt. Nam superiora illa agmina
ab intimis nunquam recedunt} (\emph{Hom.\ in Ev.} 34.12). The lesser hosts
are sent out, visibly or invisibly, for man's comfort, as angels or
archangels; the higher never withdraw from the inmost presence.

Aquinas first grants that the order of providence, in which the lower are
administered through the higher, is sometimes set aside in bodily things
for a higher order: the man born blind was enlightened and Lazarus raised
immediately by God, and God can reveal to men without angels and to higher
angels without lower. ``From this standpoint some have said that according
to the general law the superior angels are not sent, but only the
inferior; yet that sometimes, by Divine dispensation, the superior angels
also are sent.'' Aquinas rejects the exception. The angelic order follows
the gifts of grace, and ``the order of grace has no order above itself for
the sake of which it should be passed over; as the order of nature is
passed over for the sake of grace.'' Nature is passed over in miracles to
confirm faith, but passing over the angelic order would confirm nothing,
``since this could not be perceived by us.'' And nothing in the divine
ministries exceeds the lower orders: ``those who announce the highest
things are called archangels,'' and so ``the archangel Gabriel was sent to
the Virgin Mary,'' which ``was the greatest of all the Divine ministries.''
Gregory had said so in the homily: \latin{Ad hoc quippe ministerium summum
angelum venire dignum fuerat, qui summum omnium nuntiabat} (\emph{Hom.\ in
Ev.} 34.8). ``Thus with Dionysius (Coel. Hier. xiii) we must say, without
any distinction, that the superior angels are never sent to the external
ministry.''

The replies take up Scripture's apparent exceptions. Hebrews calls all the
angels ministering spirits. As the divine Persons have a visible and an
invisible mission, so there is an external mission regarding bodies, ``and
on such a mission not all the angels are sent,'' and an interior mission
regarding some intellectual effect, as when one angel enlightens another,
``and in this way all the angels are sent''; or the Apostle may mean only
the angels through whom the Law was given (ad~1). Isaiah's seraph is the
chief difficulty (ad~2). ``One of the seraphims flew to me, and in his hand
was a live coal \dots\ And he touched my mouth, and said: Behold this hath
touched thy lips, and thy iniquities shall be taken away'' (Isa 6:6--7).
Dionysius gives two answers in the thirteenth chapter of the
\emph{Celestial Hierarchy} (\S\ref{sec:ch13}), and Aquinas takes both.
First, ``the angel who was sent to purify the prophet's lips was one of
the inferior order; but was called a `Seraph,' that is, `kindling' in an
equivocal sense, because he came to `kindle' the lips of the prophet.'' In
Parker's rendering, some say that ``some one of the Angels, placed over us
as a sacred Minister of the Prophet's cleansing,'' bears the name of the
Seraphim ``on the ground that the removal of the faults spoken of, and the
restoration of him who was cleansed for the Divine mission, was through
fire'' (\emph{CH} 13.2). Gregory gives the same rule: \latin{ii
spiritus qui mittuntur eorum vocabulum percipiunt quorum officium gerunt},
and the angel who brings the coal \latin{seraphim vocatur, quod incendium
dicitur} (\emph{Hom.\ in Ev.} 34.12). Second, ``the superior angels
communicate their own proper gifts whereby they are denominated, through
the ministry of the inferior angels,'' as ``the Pope is said to absolve a
man when he gives absolution by means of someone else.'' This is the answer
Dionysius received from ``another man'' whose defence was ``by no means
foolish'': the angel ``referred his own cleansing function to God, and after
God, to the first working Hierarchy'' (\emph{CH} 13.3), and so ``attributes
his own purifying science and power to God, indeed, as Cause, but to the
Seraphim as first-operating Hierarch'' (13.4). Dionysius's own comparison
is the bishop who purifies through his priests. The divine Persons are
sent only in an equivocal sense, not in ministry (ad~3). Among unequal
angels there is a scale of ministries, so that the higher among the sent
are sent to the higher ministries and the lower to the lower (ad~4).

Scripture and the Fathers show the lower angels sending one another. ``And
behold the angel that spoke in me went forth, and another angel went out
to meet him. And he said to him: Run, speak to this young man'' (Zach
2:3--4). Gregory concludes from it, \latin{Dum enim angelus ad angelum
dicit Curre et loquere, dubium non est quia alius alium mittit. Minora vero
sunt quae mittuntur, majora quae mittunt} (\emph{Hom.\ in Ev.} 34.13), and
Dionysius reads the same vision of one of the first angels instructing an
angel of lower rank (\emph{CH} 8.2).

\paragraph{The sent also assist, in one sense (a.~3).} The \emph{sed contra}
is Gregory on Job 25:3, ``Is there any numbering of his soldiers?'':
``Those powers assist, who do not go forth as messengers to men''
(\emph{Moralia} XVII). At the locus Gregory reads the verse with Dan 7:10.
``The number of the citizens above is represented as infinite and definite,
in order that that which relatively to God is capable of being numbered
may be shewn relatively to man to be incapable of being numbered.'' Then he
distinguishes: ``those Powers stand before Him without a doubt, which never
go forth for the communicating things to men,'' while those who minister
come as bearers of tidings, ``yet these same beings also, by the act of
contemplation, are not withdrawn from the interior world'' (\emph{Moralia}
XVII.13.18). The homily says the same: \latin{Aliud namque est ministrare,
aliud assistere, quia hi administrant Deo, qui et ad nos nuntiando exeunt;
assistunt vero qui sic contemplatione intima perfruuntur, ut ad explenda
foras opera minime mittantur} (\emph{Hom.\ in Ev.} 34.12). Gregory also
gives the angels their name as soldiers: ``we are not unaware that those
war against the powers of the air, which same conflicts however they carry
on not by labour but by authority'' (\emph{Moralia} XVII.13.19).

Aquinas explains the two words ``after the likeness of those who attend
upon a king; some of whom ever wait upon him, and hear his commands
immediately; while others there are to whom the royal commands are
conveyed by those who are in attendance---for instance, those who are
placed at the head of the administration of various cities; these are said
to administer, not to assist.'' ``All the angels gaze upon the Divine
Essence immediately,'' and in that respect all, even those who minister,
assist, as Gregory says that those sent on the external ministry of our
salvation ``can always assist and see the face of the Father'' (\emph{Mor.}
II). ``Yet not all the angels can perceive the secrets of the Divine
mysteries in the clearness itself of the Divine Essence; but only the
superior angels who announce them to the inferior,'' and in that respect
only the first hierarchy is said to assist, ``whose special prerogative it
is to be enlightened immediately by God.'' The first two objections, from
Gregory's ``the angels are sent, and assist'' and from Raphael, ``one of
the seven, who stand before the Lord'' (Tob 12:15), speak of the first
mode of assisting. Satan in Job was not among the assistants but present
among them (Job 1:6); ``as Gregory says (Moral. ii), `though he has lost
beatitude, still he has retained a nature like to the angels'\,'' (ad~3).
At its locus Gregory adds the image that answers how Satan could stand
before God without seeing him: ``he is said to have come before the Lord,
but not that he saw the Lord. For he came to be seen, and not to see \dots\
as when a blind man stands in the sun, he is himself bathed indeed in the
rays of light, yet he sees nothing of the light, by which he is
brightened'' (\emph{Moralia} II.4.5). All the assistants see some things
immediately in God's glory, the higher more than the lower, ``as also among
those who assist the king, one knows more of the king's secrets than
another'' (ad~4).

\paragraph{Which orders are sent (a.~4).} The article asks whether all the
angels of the second hierarchy are sent. The \emph{sed contra} is
Dionysius: ``the `Dominations are above all subjection.' But to be sent
implies subjection.'' Parker renders the eighth chapter: the name of the
Lordships ``denotes a certain unslavish elevation, free from all
grovelling subserviency \dots\ superior to every kind of cringing slavery,
indomitable to every subserviency'' (\emph{CH} 8.1). Aquinas argues from
the names, since ``the angelic properties are manifested by their names''
(\emph{CH} 7). To be sent to outward ministry is to act by divine command
on some bodily creature, the execution of a divine ministry. ``But the name
`dominations' does not signify any such administration, but only
disposition and command in administering. On the other hand, the names of
the inferior orders imply administration, for the `Angels' and
`Archangels' are so called from `announcing'; the `Virtues' and `Powers' are
so called in respect of some act; and it is right that the `Prince,'
according to what Gregory says (Hom. xxxiv in Evang.), `be first among the
workers.' Hence it belongs to these five orders to be sent to external
ministry; not to the four superior orders.'' The Dominations are counted
among the ministering angels as the architect is among the builders,
``disposing and commanding what is to be done by others'' (ad~1).

The second reply records a real difference between Gregory and Dionysius
over Dan 7:10. Gregory reads ``thousands of thousands ministered to him''
partitively, thousands out of thousands, an indefinite number signifying
excess, and ``ten thousand times a hundred thousand stood before him'' as a
finite number; so the ministers are more than the assistants. This
``rests on the opinion of the Platonists, who said that the nearer things
are to the one first principle, the smaller they are in number,'' and it
``is verified as regards the number of orders, as six administer and three
assist.'' Dionysius holds instead that the multitude of angels surpasses
all material multitude, ``because whatever is better is more intended and
more multiplied by God. Hence, as the assistants are superior to the
ministers there will be more assistants than ministers,'' and the verse is
read by multiplication, a thousand times a thousand, ``ten thousand times a
hundred thousand'' being far more. The figure does not give the number of
the angels but signifies that it ``exceeds all material multitude'' by
multiplying the greatest numbers, ``as Dionysius remarks in the same
passage'' (ad~2; \emph{CH} 14). Gregory's own words at the locus are that
``they are more in number that `minister' than those that preeminently
`stand before Him''' (\emph{Moralia} XVII.13.18), and Dionysius says that
the hosts of the supermundane minds are ``surpassing the weak and
contracted measurement of our material number'' (\emph{CH} 14).

The offices by which the orders serve the government of the world, as the
\emph{Contra gentiles} assigns them, stand beside the answer of q.~112
a.~4 on their outward mission. On the two orderings of the orders see
Appendix~\ref{app:orders} and \S\ref{sec:q108}.

\begin{threestable}{Order}{Office in the government of the world
(\emph{SCG} III.80)}{Outward mission (\ST{I q.~112 a.~4})}
Seraphim & See the reason of providence in the last end, the divine
goodness, whence their name of burning (n.~5) & Not sent\\
Cherubim & See it in the divine form, the fullness of knowledge (n.~6) &
Not sent\\
Thrones & Consider the disposition of the divine judgments in itself
(n.~7) & Not sent\\
Dominations & Command what others execute, \latin{dominorum enim est
praecipere} (n.~10) & Not sent; they dispose and command\\
Virtues & Move the heavens and carry out works beyond the order of nature
(n.~11) & Sent\\
Powers & Guard the order of providence against what would disturb it
(n.~12) & Sent\\
Principalities & Rule the common good of nations and kingdoms (n.~14;
Dan 10) & Sent\\
Archangels & Announce what all must believe, as Gabriel the Incarnation
(n.~15) & Sent\\
Angels & Serve the good of each person; hence the guardians of men
(n.~16; Ps 90[91]:11) & Sent\\
\end{threestable}

\angeldossier{\ST{I q.~112 aa.~1--4}; parallels \emph{SCG} III.79--80.}
 {Common teaching, from Scripture, that angels are sent by God to
 minister to men (a.~1; Exod 23:20; Heb 1:14; Luke 1:26), and that those
 sent never lose the vision and contemplation of God (a.~1 ad~3, a.~3;
 Matt 18:10; Gregory, \emph{Moralia} II.3.3, \emph{Hom.\ in Ev.} 34.13;
 Augustine, \emph{De Gen.\ ad litt.} VIII.24.45). Thomist position, with Dionysius and Gregory, that the
 superior angels are never sent to outward ministry and that the lower
 bear the names of the higher whose gifts they carry (a.~2; \emph{CH} 13;
 \emph{Hom.\ in Ev.} 34.12), and that the four highest orders are not sent
 while the five lower are (a.~4). Common teaching, from Dan 7:10 as
 Gregory reads it, that some angels assist and others minister; Thomist
 position that all assist by the vision of God and the first hierarchy
 alone by immediate illumination in the divine mysteries (a.~3).
 Disputed between Gregory and Dionysius whether the ministering or the
 assisting angels are the more numerous (a.~4 ad~2).}
 {Exod 23:20 (\emph{sed contra} of a.~1); Gregory, \emph{Hom.\ in Ev.}
 34.8, 34.12--13 (\emph{sed contra} of a.~2; aa.~2--4), \emph{Moralia}
 II.3.3, II.4.4--5 (aa.~1, 3) and XVII.13.18--19 (\emph{sed contra} of
 a.~3; a.~4 obj.~2, ad~2), read at the locus; Dionysius, \emph{CH} 13.1--4
 (a.~2), \emph{CH} 8.1--2 (a.~4 obj.~1 and \emph{sed contra}), \emph{CH} 7
 (aa.~3--4), \emph{CH} 14 (a.~4 ad~2), read in Parker; Damascene,
 \emph{De fide orth.} II.3, read at the locus; Augustine, \emph{De Trin.}
 IV.20 (a.~1 ad~1), cited by Aquinas, and \emph{De Gen.\ ad litt.}
 VIII.24.45, read at the locus; Heb 1:14, Isa 6:6--7, Tob 12:15, Job 1:6,
 Job 25:3, Dan 7:10, Zach 2:3--4, Luke 1:26, Luke 22:27, Phil 2:3, Ecclus
 38:25.}
 {Those who held that the superior angels are sometimes sent by divine
 dispensation (a.~2). Gregory's partitive reading of Dan 7:10, on which the
 ministers outnumber the assistants, against Dionysius's reading, on which
 the assistants are the more numerous (a.~4 ad~2). The reading of Heb 1:14
 by which all angels without distinction are sent to outward ministry,
 which Aquinas restricts to interior mission or to the angels of the Law
 (a.~2 ad~1).}
